% Volume III: Attention and Publicity
% Part VI. Privacy as Boundary and Capacity
% Chapter 29: Privacy Beyond Secrecy
% Spine: proof-spines/ch029-spine.md

\chapter{Privacy Beyond Secrecy}
\label{ch:ch029}

Privacy is defined here not as concealment but as control over the conditions under which information about a person crosses contextual boundaries. Those conditions govern collection, inference, combination, retention, retrieval, interpretation, and action. The chapter distinguishes informational privacy, decisional privacy, spatial privacy, relational privacy, inferential privacy, and temporal privacy, and treats each as a mode of boundary maintenance rather than a mere reservoir of hidden facts.

This broader account explains why the slogan ``nothing to hide'' misses the point. Privacy is a capacity for agency, rehearsal, role differentiation, and selective disclosure under non-totalizing conditions. It is therefore compatible with, and often necessary for, responsible inquiry: a society can demand accountability from power without demanding exhaustive legibility from persons.
